\chapter{General description}\label{ch:description}
%----------------------------------------------------------------------
% This chapter describes the general factors that affect the product
% and its requirements. The chapter  does not state specific
% requirements but makes those requirements easier to understand.

%----------------------------------------------------------------------
\section{Product perspective}
%----------------------------------------------------------------------
% This section puts the product into perspective with other related
% systems. If the product is to replace and existing system, the
% system should be described and referenced Ancestors of the product
% that are no longer in use might be mentioned. If the product is 
% ``stand alone'', it should be stated here.

The \ssat{} product will be based on the following existing products:

\begin{itemize}
\item \emph{GNAT Pro}, a commercial Ada compilation system and product
  support.
\item \emph{GPS}, the GNAT Programming Studio.
\item \emph{GCC}, the GNU Compiler Collection.
\end{itemize}

The resulting product will be part of the GNAT Pro family, which
includes Ada compilation systems and related tools for a large set of
host and target platforms.

%----------------------------------------------------------------------
\section{General Capabilities}
%----------------------------------------------------------------------
% This section should describe the main capabilities and why they are
% needed. This section should describe the process to be supported by
% the software, indicating those parts of the process where its is used.

The \ssat{} to be developed is a commercially supported product that
will anticipate the stack space requirements of applications and allow
safe prevention of overflow conditions.

The Static Stack Analysis Tool will analyze the application stack
space consumption patterns and compute worst-case bounds prior to
execution time. It will be done by performing local stack consumption
analysis combined with control-flow graph traversals.

%----------------------------------------------------------------------
\section{General Constraints}
%----------------------------------------------------------------------
% This section should describe any item that will limit the developer's
% options for building the software.  This section should not be used to
% impose specific requirements or specific design constraints, but
% should state the reasons why certain requirements or constraints exist.

The tool will use the compiler for extracting per-function stack usage
and per-unit call graph information. Hence, \ssat{} will only be able to
analyse stack consumption on units that can be compiled (it excludes
third-party libraries not providing the required information), and the
results will only be valid for the final application compiled with the
same compiler (and the same compilation options).

%----------------------------------------------------------------------
\section{User characteristics}
%----------------------------------------------------------------------
% This section should describe those general characteristics of the
% users affecting the specific requirements.  Many people may interact
% with the software during the operations and maintenance phase.  Some
% of these people are users, operators and maintenance
% personnel. Certain characteristics of these people, such as
% educational level, language, experience and technical expertise impose
% important constraints on the software.  Software may be frequently
% used, but individuals may use it only occasionally. Frequent users
% will become experts whereas infrequent users may remain relative
% novices. It is important to classify the users and estimate the likely
% numbers in each category.

The expected users of the \ssat{} are software engineers working on Ada
and/or C projects. The users should have working knowledge of the
English language and technical expertise in software development in
Ada and/or C.

%----------------------------------------------------------------------
\section{Operational environment}
%----------------------------------------------------------------------
% This section should describe the real world the software is to operate
% in. Context diagrams may support this narrative description, to
% summarise external interfaces and system block diagrams, to show how
% the activity fits within the larger system. The nature of the
% exchanges with external systems should be specified.  If a System
% Specification defines a product that is a component of a parent system
% or project, then this section should outline the activities that will
% be supported by external systems References to the Interface Control
% Documents that define the external interfaces with the other systems
% will be provided. The computer infrastructure to be used will be also
% described.

The \ssat{} will operate on any system on which the \emph{GNAT Pro}
compiler is supported, including both native and cross compilers. The
particular versions of the operating systems on which the software is
to be validated will be defined within software requirements analysis
process.
